%%%%%
%% DIESES DATEI IST ALT UND WIRD NICHT MEHR GEBRAUCHT
%%%%%





\chapter{Installation Guide}
\label{cha:installation guide}

\section{System Requirements}

Intel-PC or Sparc-compliant machine running Solaris 2.3 or higher.
32 MB main memory and approximately 20 MB disk space.
The Solaris Binary Compatibility package must be installed.
You can check this with the 'pkginfo' command.
A recent Sun Recommended Patch Cluster should be installed.

If you receive errors regarding insufficient swap space you
must configure more; for example as root do a:

\begin{verbatim}
mkfile 100m /export/home/further_swap
swap -a /export/home/further_swap
\end{verbatim}


Make this swap permanent by appending the following line
to /etc/vfstab:

\begin{verbatim}
/export/home/further_swap - - swap - no -
\end{verbatim}

In order to uncompress the CB tar file you need the
\verb+gunzip+ utility which is available for free on many ftp-servers.


\section{Files}

\begin{description}
\item[cb\relnrwodot .tar.gz:] ConceptBase executables (sun4/i86pc), documentation and scripts
\item[CBinstall:] Script to install ConceptBase on your system
\item[InstallationGuide:]   this section in plain ASCII
\end{description}

\section{Installation}

To obtain the system, please transfer \verb+CBinstall+ and {\tt cb\relnrwodot .tar.gz}
 to your local machine.
To install the system, ensure that the script \verb+CBinstall+ is executable.
If not, change the permission mode and run \verb+CBinstall <directory>+,
where \verb+<directory>+ is the directory where ConceptBase will be installed.

The script will install the system and create the file \verb+login.new+,
which contains the settings of all necessary environment variables.


The graphbrowser tool needs special Andrew-Toolkit fonts to
operate correctly. These fonts must be accessible through
the X11-fontpath. To ensure this, extend the fontpath with
the following command:

\begin{verbatim}
xset fp+ $CB_HOME/lib/andrew6.3/X11fonts
\end{verbatim}

We recommend to

\begin{itemize}
\item create a special user \verb+cbase+ to be responsible for the installation
  and owning all installed files.

\item create a group \verb+cbase+ in /etc/groups and assign each
  ConceptBase user to it
\end{itemize}


\section{Registering your installation}

Before ConceptBase can be started, you have to decode it.
To do this, please perform the following steps:

Register your installation via the URL

\begin{latexonly}
\verb+http://www-i5.informatik.rwth-aachen.de/CBdoc/onlineReg.html+
\end{latexonly}
\html{\htmladdnormallink{Online Registration}{http://www-i5.informatik.rwth-aachen.de/CBdoc/onlineReg.html}}

or via sending an email to

\begin{latexonly}
\verb+cb@I5.informatik.rwth-aachen.de+
\end{latexonly}
\html{\htmladdnormallink{ConceptBase Team}{mailto:cb@I5.informatik.rwth-aachen.de}}

This email must contain the full postal adress of your institution.

You'll receive an email, which contains decode information.
This email must be saved into a file named \verb+regInfo+, that has to
be placed into the directory

\begin{verbatim}
$CB_HOME/lib/system
\end{verbatim}

Read the license agreement, placed in
\verb+$CB_HOME/doc+, sign it and send it to

\begin{verbatim}
ConceptBase-Team
c/o Christoph Quix
RWTH Aachen - Informatik V
D-52056 Aachen - Germany
\end{verbatim}



\section{Starting ConceptBase}

\subsection{Starting the ConceptBase Server}

You have to use a C-Shell (/bin/csh) to start ConceptBase. Start the server by

\begin{verbatim}
$CB_HOME/bin/CBserver -p 4001 -d DB
\end{verbatim}

The number 4001 is the port (or {\em socket}) number that is used to connect
 clients to the server process. You can specify any number between 2000 and
99999. The default value is 4001 (if you omit the '-p' option). The parameter
-d specifies the database (sometimes called {\em application} in the user manual)
to be loaded by the server at startup time. If a path is omitted (as above)
then the database is loaded from the current directory. If there is no
database with the specified name (here:DB) then it is created and initialized
with the system classes of O-Telos, the object model of ConceptBase.
After successfully loading the server reports

\begin{verbatim}
          This is ConceptBase (CBserver) ...
\end{verbatim}

with some additional information.

Please note that no two servers can use the same port number on the same
machine. In this case you get the error message

\begin{verbatim}
IPC Error: Unable to bind socket to name
\end{verbatim}

and the server won't start up. You should then either start the server on a
different machine or use another port number.

The list of other command line parameters can be obtained by
typing \verb+CBserver -h+. Specified parameters must be separated by blanks.



\subsection{Starting the ConceptBase Usage Environment}

Start the X11 usage environment of ConceptBase in a C-Shell by

\begin{verbatim}
$CB_HOME/bin/CBinterface
\end{verbatim}

Successfully starting is indicated by the appearing CBworkbench.

If you have installed the Java Development Kit or Java Runtime
Environment (1.1 or newer) you can start the ConceptBase Java
Interface by

\begin{verbatim}
$CB_HOME/bin/CBjavaInteface
\end{verbatim}

\subsection{Connecting Server and Usage Environment}

After starting server and usage environment of CB the first step is to
establish the connection between them. This is done by selecting menu item
{\em Connect CB server}
from the Server  menu of CBworkbench. In the appearing interaction window the
host name and port number of the CBserver must be specified and as a result
the connection is established.

\begin{flushleft}
{\bf Example:}
\end{flushleft}
\begin{enumerate}
\item start server on machine alpha with port number 5000 (see above)
\item start usage environment on machine beta (see above)
\item connect usage environment on beta to the server on alpha
              	by selecting Connect CB server and specifying
				as hostname \verb+alpha+ and port number 5000.
\end{enumerate}

Whenever a new connection from a usage environment to a CBserver is created a
possibly existing 'old' one is closed. Please note that the machines
alpha and beta can be located anywhere in the Internet. You can also start
both processes on the same machine.



\section{Important additional remarks}

A CBserver may only be stopped by a client if the name of the user of the
client is the same as the one who started the CBserver. Use menu item Report
clients of the usage environment to display the user names. This behavior is
not an error but a security measure.

Objectbases created with previous releases of ConceptBase (i.e. 4.x and earlier) are not compatible
with release 5.0 or later, due to the completely new storage management system.
To cope with this you have the possibility to bulk extract frames
(via \verb+Store Frames+) from the previous server and feed them into the new server.

We hope the system is useful for your application. We recommend to start with
the (in
\htmladdnormallink{ConceptBase Tutorial}{../../Tutorial/html/cbTutorial.html}
\begin{latexonly}
\verb+$CB_HOME/doc/Tutorial+
\end{latexonly}
) for getting used to the language and the system.

A bibliography on the system is available via the following URL:\\
\begin{latexonly}
\verb+http://www-i5.informatik.rwth-aachen.de/CBdoc/cblit.html+
\end{latexonly}
\html{\htmladdnormallink{Bibliograhpy}{http://www-i5.informatik.rwth-aachen.de/CBdoc/cblit.html}}



\section{Contact}

\begin{verbatim}
ConceptBase-Team
RWTH Aachen - Informatik V
D-52056 Aachen - Germany

Tel./Fax: +49-241 80 21 501 / +49-241-80 22 321
\end{verbatim}
\begin{latexonly}
\verb+Email: cb@I5.informatik.rwth-aachen.de+ \\
\verb+WWW: http://www-i5.informatik.rwth-aachen.de/CBdoc/cbflyer.html+ \\
\end{latexonly}
\html{\htmladdnormallink{Email}{mailto:cb@I5.informatik.rwth-aachen.de}}
\html{\htmladdnormallink{ConceptBase WWW Site}{http://www-i5.informatik.rwth-aachen.de/CBdoc/cbflyer.html}}

